CIĄG DALSZY ROZDZIAŁU 14 (ze strony 9-10):

długo, jak jest przez nią noszony.

Rozpoznawanie mikstur – Medalion pozwala rozpoznać rodzaj mikstury bez próbowania jej (picia).
Tak więc znaleziona mikstura zachowuje się w całości i może być w przyszłości użyta dwa razy.

Elokwencja – Pozwala na dodanie 3 punktów do liczby oczek, wyrzuconych podczas ustalania wyniku negocjacji.
Postać, nosząca medalion nie musi być tą, która prowadzi negocjacje – wystarczy, że znajduje się ona w tym samym segmencie.

Zręczność – Medalion podnosi zręczność postaci, która go nosi o 2 punkty.

Dusiciel – Medalion zaciska się na szyi postaci, która go założy, zadając jej 1K6 ran.

[14.8] Pierścień.

Odporność +1 lub +2 – Podnosi odporność postaci noszącej ten pierścień odpowiednio o 1 lub 2 punkty.

pozostałe – Pozwalają postaci, noszącej jeden z nich na rzucenie za darmo (bez ponoszenia kosztu) czaru, odpowiadającego nazwie pierścienia, ale tylko jeden raz w ciągu przygody.

[14.9] Tabele ,,Kosztowności'', ,,Magiczne przedmioty'', ,,Broń'' – patrz plansza z tabelami.

15.0 SPECJALNE MOŻLIWOŚCI NIEKTÓRYCH POTWORÓW

Ogólna zasada:

Niektóre potwory posiadają specjalne możliwości, które są wykorzystywane podczas walki z oddziałem.

Zasady szczegółowe:

[15.1] Kronki śmierdzą w sposób tak straszny, iż może to przyprawić o chorobę.
Jeżeli oddział napotka kronki i rozpoczyna z nimi walkę każdy z jego członków musi podjąć próbę odparcia czaru.
Postać, która tej próby nie przejdzie w sposób pomyślny zaczyna chorować i jej zręczność jest obniżona w czasie tej walki o 1 i 2 punkty.
Po walce z kronkami postać zdrowiej.

[15.2] Chimera, atakując postać w normalny sposób, może dodatkowo zionąć ogniem.
Zaatakowana postać musi próbować odeprzeć czar, a jeżeli będzie nieudana otrzymuje, poza ranami, odniesionym na skutek normalnego ataku chimery, jedną dodatkową ranę od ognia.

[15.3] Meduza może zamienić zaatakowaną postać w kamień.
Przy każdym spotkaniu z Meduzą należy rzucać 1K6.
Jeżeli zostanie wyrzucone 6 oczek zaatakowane przez Meduzę postać zamienia się w kamień.
Żeton, reprezentujący tę postać powinien zostać usunięty z szyku oddziału.

Jeżeli któryś z członków oddziału zna czar ,,Ożywienie kamienia'', może, bezpośrednio po zakończeniu walki, przywrócić zamienioną w kamień postać do życia.
Jeżeli nikt nie zna tego czaru postać jest uważana za martwą.

[15.4] Troll posiada zdolność regeneracji.
Po zakończeniu każdej trzeciej fazy ,,atak potwora'' Troll zmniejsza o jeden liczbę ran, które odniósł.

[15.5] Wampir posiada zdolność rzucania ,,za darmo'' czaru ,,urok'' ale tylko podczas dwóch pierwszych rund walki.
Jeżeli zaatakowanej tym czarem postaci nie uda się go odeprzeć dostaje się pod władzę Wampira i jest umieszczana w pierwszym szeregu szyku potworów.
Czar może być złamany tylko przez śmierć Wampira lub przez czar ,,uwolnienie''.
Postać, znajdująca się pod wpływem uroku jednak musi w inny sposób atakować oddział.

[15.6] Każda rana zadana Hydrze zwiększa jej zręczność o jeden punkt.
Tak więc w toku walki zdolności bojowe Hydry rosną wprost proporcjonalnie do liczby odniesionych ran.
Hydra jest uśmiercana na normalny sposób, tzn. jeżeli liczba zadanych jej ran jest równa się ze współczynnikiem siła.

[15.7] Czarnosiężnicy (wraz z Bezimiennym) znają jeden czar – ,,błyskawica''.
Używają go w trakcie każdej fazy ,,atak potwora''.
Czarnoksiężnik nie może użyć tego czaru, jeżeli związany z nim koszt spowodowałby jego śmierć.